\section{Discussion}
\label{sec:discussion}

\subsection{What text supervision buys, and where it stops}

We find that performance tracks organ size and contrast almost monotonically
(Table~\ref{tab:perclass}), which is the shape expected for a method with no
spatial supervision. The prompt must select a region by appearance and position
alone, and both cues degrade together for small, low-contrast structures.

The pancreas result is the sharpest statement of the ceiling, because
it is invariant to the one thing we can freely change. We measure pancreas IoU
at 0.213--0.218 across \emph{all five} of our prompt formulations: the bare word
``pancreas,'' a generic clinical caption, position-only text, the full
multi-sentence anatomical description, and the adversarial variant. A manipulation that moves the other classes by 60 percentage points moves the
pancreas by half a point. We conclude that whatever limits the pancreas here is
not in the text, so prompt engineering will not repair it. The constraint we run
into is that a small, low-contrast, elongated structure does not present a
$14\times14$-patch signature the encoder can separate from surrounding soft
tissue.

The distinction between prompt-limited and data-limited classes is worth drawing
explicitly, because the two call for different responses. The kidneys sit at 0.45--0.49 IoU and collapse by 47--63\% under adversarial
position swaps, which tells us they depend heavily on the positional text and
are therefore prompt-limited: better laterality cues should help them. The pancreas does not move under any prompt manipulation at all, so no amount of
prompt engineering will recover it, and the productive direction is a finer
patch grid or a different backbone. Reading only the aggregate mIoU would hide
that difference and invite us to spend effort in the wrong place.

\subsection{What the patch objective actually contributed}

The largest single gain comes from adding the patch-level loss, which moves mIoU
from 0.063 to 0.307. The size of that step is the most informative
number in Table~\ref{tab:progression}, because it quantifies how little
localization a globally-aligned model gives up for free. Before the patch loss
we were fine-tuning the same backbone with the same data and the same text
encoder, differing only in whether any objective required individual patch
tokens to be discriminative. Without such an objective, the patch tokens of a
contrastively trained encoder are optimized only insofar as they contribute to a
pooled embedding, and a pooled embedding can be correct while its constituents
are individually uninformative.

This is consistent with what the dense-alignment literature reports from the
opposite direction. MaskCLIP finds that CLIP retains more spatial information
than the pooled representation reveals, and recovers it by bypassing pooling
\cite{zhou2022maskclip}, and later work extends that recovery to weakly
supervised semantic segmentation \cite{yang2025clipdense}. Analyses of CLIP's
explainability reach a similar conclusion about why raw attention localizes
poorly \cite{li2024clipexplain}. We read our 0.063 measurement as the same
phenomenon observed without any recovery step: the information is latent rather
than absent, and something must either extract it or train it in.

We chose to train it in, and the reason generalizes. Extraction methods
recover \emph{where} an organ is, which is sufficient for segmentation
benchmarks. The downstream index additionally needs organ maps that vanish
where the organ is not present, because a proximity score computed against a
phantom mask is meaningless. That is a calibration property rather than a
localization one, and we found it easier to obtain from an objective that
explicitly pushes non-organ patches away than from a post-hoc reading of a model
never trained to express absence.

\subsection{Absence is information, and normalization destroys it}

The calibration correction added 0.059 mIoU and halved HD95 without retraining,
which makes it the cheapest improvement we report. The underlying lesson is more transferable than the number.

Per-image min--max normalization is close to a reflex when visualizing saliency,
and it is actively harmful here. Rescaling every slice to a fixed range forces a
slice containing no pancreas at all to span the same output range as a slice
centred on the pancreas. The relative ordering within each slice survives, but
the comparison \emph{between} slices, and therefore any notion of confidence,
does not. The raw similarities are bounded in $[-1,1]$ and routinely negative
across an entire slice, and that negativity is signal rather than noise: it is how the model expresses that the prompted organ is not here.

Any pipeline that reads similarity maps from a vision-language model should
check whether it is discarding this. The failure is silent, because
normalized maps look better in figures and score similarly on any metric
computed slice-by-slice among organ-present slices. It surfaces only when the
maps are asked a question that requires absolute comparison across slices, which
is exactly what the 3D index does.

\subsection{Elaboration is not what makes prompts work}

The caption ablation licenses a stronger conclusion than ``text matters.'' We
find the class name alone recovers 96.9\% of the full description's mIoU, and we
find the remaining gap is closed by a one-line generic clinical caption carrying
no anatomical content. L2 and L4 differ by 0.0006 on both metrics and split
them, with L2 ahead on mIoU and L4 ahead on Dice. The elaborate multi-sentence
descriptions we used throughout training therefore bought us nothing measurable
over a single generic sentence.

This matters practically, because writing those descriptions was the most
labour-intensive part of the setup, and the result says that labour was largely
wasted. Anyone extending this approach to a new organ set
should write class names and one-line captions, verify that this suffices, and
invest the saved effort in the failure modes below. One qualification applies: our descriptions were used during training as well as evaluation, so
we cannot separate the possibility that the model has already absorbed the
anatomical content from the possibility that the content was never necessary.
Distinguishing those would require retraining with short captions, which we did
not do.

The L3 result sharpens what the positional language is for. Position without
identity is actively worse than identity alone (0.561 vs.\ 0.607), so we
conclude the position clauses are not carrying the localization. They refine a
target the class name has already selected.

\subsection{Failure is directional, which is the safety-relevant finding}

The accidental error in the L3 \org{lung-right} prompt is the most operationally
important result in this paper, and we did not design it. The prompt named the wrong organ \emph{and} the wrong side, consistently, so it
became a valid description of the left lung. The model did not degrade, hedge,
or produce a diffuse map. It confidently segmented the left lung, and we measured
IoU on \org{lung-right} falling to 0.23 while \org{lung-left} at the same prompt
level held at 0.829, which proves the level itself was harmless and isolates the
cause to two words.

We therefore observe that a wrong-but-coherent prompt yields a confident
segmentation of the wrong structure. For any deployed system this is a worse
failure mode than low confidence, because nothing in the output signals
that anything went wrong. The mask is clean, well-formed, and anatomically
plausible. It is simply of a different organ. We argue that systems built on
text-prompted localization need prompt validation as a first-class component,
not an afterthought, because a wrong prompt does not announce itself downstream.

This failure mode is a direct consequence of the property the paper spends its
length establishing. A model that genuinely follows its prompt will follow a
wrong prompt just as faithfully as a right one, and the same evidence that
proves the text is load-bearing (finding 3) also proves that corrupting the text
is dangerous. This is not a defect specific to our pipeline. Any
method in which the prompt truly determines the output inherits it, and methods
that place a trained refiner between prompt and mask may partially mask the
symptom without removing the cause.

Finding 3a gives a concrete mitigation. We saw \org{lung-right} survive
the adversarial swap far better than \org{lung-left} ($-19.5\%$ vs.\ $-56.8\%$)
because its prompt contains a clause the swap cannot corrupt: the right lung
genuinely has three lobes and the left has two, so ``three lobes are present''
still points at the correct organ after every laterality word is flipped. We
conclude that prompt robustness is controllable by construction. If we include
morphological cues that are invariant under position swaps, such as lobe count,
characteristic shape, or vessel pattern, a bilateral class stops depending
entirely on laterality words it cannot verify. The swap test itself is cheap enough to be a routine design check: it
requires no retraining, and a class that collapses under it is a class whose
prompt rests entirely on assertions the image cannot confirm.

\subsection{Where our approach sits among dense-alignment methods}

We do not claim the best boundaries. DenseCLIP fine-tunes with a dense decoder
\cite{rao2022denseclip}, CLIPSeg attaches a transformer decoder emitting a
segmentation from a prompt \cite{lueddecke2022clipseg}, and MedCLIP-SAMv2 hands
mask formation to SAM \cite{koleilat2025medclipsamv2, kirillov2023sam}. All of
these will produce cleaner edges than an upsampled $14\times14$ grid.

What we gain by not doing this is attribution. Our caption ablation is
interpretable precisely because nothing sits between the prompt and the mask: a
55-point IoU drop can be traced to two words, because there is no intermediate
component that might have absorbed or amplified the error. Had a trained
decoder been in the path, we could not have said whether the decoder repaired a
bad localization or a good one was corrupted downstream. The architectural simplicity is the enabling
condition for the prompt findings rather than an oversight, and any attempt to
replicate those findings on a refiner-based pipeline would face this confound.

GLoRIA is the closest medical antecedent, learning global and local
representations jointly by matching sub-regions to report words
\cite{huang2021gloria}. The important difference is the supervision target. GLoRIA attends over words to build localized representations for
label-efficient recognition, whereas we supervise patch tokens directly against
a mask so the similarity map is the output rather than an intermediate. That
choice is what makes the map usable as geometry in our index, and it is also
what makes our failures legible as prompt failures.

\subsection{Our failure band is narrower than it first appears}

Describing T12--L2 and L5 as a uniform ``thoracolumbar failure'' overstates the
evidence, because those four levels fail in three different ways
(Table~\ref{tab:perlevel}).

\textbf{T12 and L5 are at chance, not below it.} We measure T12 missing by 0.10
rank (1.60 vs.\ 1.50) and L5 by 0.08 (3.08 vs.\ 3.00). T12 in particular faces the most stringent bar in the study: with seven expected organs its chance mean
rank is 1.50, so our model must rank a correct organ essentially first, every
time, to register any margin at all. Missing that bar by a tenth of a rank is not
comparable to failing a lenient one.

\textbf{L1 is a modest miss against a stringent bar}, at $+0.41$ against chance
2.00 with $k=5$.

\textbf{L2 is our only genuine failure}, at $+0.90$ against a lenient chance of
3.00 with only three expected organs. It is the one level where the model is
clearly worse than guessing, and it deserves to be reported as such
rather than averaged into a band.

\subsection{Two classes account for much of our damage}

The per-lesion detail identifies a more specific culprit than ``the lumbar spine
is hard.'' Across all 180 lesions we count the following.

\begin{itemize}
\item \org{bladder} takes top rank on 29 lesions, of which \textbf{21 are
      anatomically wrong and 8 correct}. It wins nearly three times more often
      where it does not belong than where it does.
\item \org{femur-left} and \org{femur-right} take top rank on 14 lesions, of
      which \textbf{zero are correct}. Our femur classes never once rank first
      at a level where a femur is anatomically expected, while winning 14 times
      where one is not.
\end{itemize}

This reframes our phantom problem. We are not looking at a generally-sound model
that occasionally misfires at difficult levels. Two of our eleven classes are
mislocalized in a way that produces false positives broadly and contributes
nothing where they would be correct. The two failing levels follow directly. At
L2, \org{bladder} hijacks top rank on 4 of 10 lesions at a level in the upper
lumbar spine where the bladder is nowhere near. At L5, the one level where
femurs \emph{are} expected, neither of our femur classes ever wins, and the
level's correct detections come from \org{bladder} instead.

It is unlikely to be coincidence that these are precisely the two classes the
text-only prompt diagnostic flagged. The \org{femur} prompt's appearance clause (``bright hyperdense cortical ring $\ldots$ central medullary canal'')
embeds almost identically to an undefined \org{humerus} prompt, more closely
than to any of femur's own defined classmates, and it describes a rib in
cross-section about as well as it describes a femur. The qualitative panel
(Figure~\ref{fig:qualitative}) confirms this visually: the winning
\org{femur-right} patches sit on rib cross-sections around the chest wall, on no
structure resembling a femoral shaft. The class is not localizing a femur at all. It is matching a generic cortical ring, and the 11-class vocabulary contains
nothing defined to outbid it.

These counts rest on small per-level samples ($n=10$ at L2, $n=13$ at L5) and
should be read as a diagnosis to act on rather than a precise effect size.

\subsection{Why both our calibration variants failed}

Both negative results have identified causes, which is what makes them worth
reporting rather than merely disclosing.

\textbf{Our global hard negatives failed by over-application, not by being the
wrong idea.} We added rib, humerus and aorta negatives to target exactly the
confusion we diagnosed above, and they worked on their intended target: femur
pair-splitting fell from 31 to 4. But we subtracted the negative from
\emph{every} organ's score rather than only from the classes it was meant to
correct, so it bled into the lungs' own scoring at the rib-lung boundary, where
lung tissue legitimately abuts rib, and created 46 new lung pair splits that had
not existed. The mechanism is specific and so is the fix it implies. We should
scope each negative to the classes we introduced it for, so that humerus and rib
compete only against our femurs and aorta only against our bladder. That per-organ-scoped variant has not
been run, and this result reads as ``wrong scoping'' rather than ``negatives do
not help.''

\textbf{Our z-score floor failed by permissiveness.} It eliminated bilateral
pair splitting entirely ($34 \to 0$) and improved mean rank against its correct
comparator ($1.94 \to 1.82$), but it more than doubled anatomically impossible
detections ($38 \to 79$). Both effects follow from one mechanism. A floor
set at $\mu + 2\sigma$ per organ is lenient enough that most organs retain a
sliver of territory almost everywhere instead of correctly vanishing where they
are absent. Slivers make bilateral territory symmetric, which zeroes our splits,
and a sliver is also enough to win a proximity contest at 0--1\,mm, which is
where nearly every new wrong detection landed.

\subsection{Evaluation-set discipline, and why we stop here}

The 180-lesion set has now absorbed several decisions: the calibration fit, the
bilateral-grouping change, the move from per-patient to per-lesion scoring, and
the two pre-registered calibration variants. It is a finite resource, and
continuing to test hypotheses against it would compound a multiple-comparisons
exposure that is already non-trivial.

This is our stated reason for not pursuing the residual phantom detections
further in this work, and we intend it to be read as a claim about what our
evaluation set can still support rather than a claim that nothing remains to be
found. The per-organ-scoped negative variant is the obvious next experiment, and the
right way to run it is against fresh data, or with
the exposure declared in advance, not as one more pass over the same 46
patients.

\subsection{Limitations of the localization method}

\textbf{Resolution.} Patch supervision is fixed at $14\times14$ by ViT-B/16 at $224\times224$.
Boundary precision is bounded by that grid, and the HD95 of 18.8\,px largely
reflects it. Raising input resolution is not a free fix: it changes the patch geometry our entire calibration and
mask-formation path is built around.

\textbf{Metric asymmetry.} All our Phase~1 numbers are 2D per-slice metrics on
organ-present slices, which is systematically more forgiving than 3D volumetric
scoring. Every cross-paper comparison here inherits that caveat, and the
supervised-to-weakly-supervised gap is wider than the tables' raw column
differences suggest.

\textbf{Prompt confound.} The anatomical descriptions were used during training
as well as at evaluation, so the caption ablation cannot fully separate ``the
description was never necessary'' from ``the model has already internalized the
description.'' Retraining with minimal captions would settle this, and we did
not do it.

\subsection{Limitations of the index}

\textbf{We use ground-truth tumor localization.} We attempt no weakly-supervised
tumor detection. We made this a scoping decision rather than an oversight:
adding a second layer of learned uncertainty before validating the index itself
would give us an uninterpretable result, since we could not attribute a failure
to the index or to the detector. It does mean our pipeline is not end-to-end
weakly supervised, and a deployed version would need that second stage.

\textbf{We fix our index parameters rather than fitting them.} We chose
$\alpha = \beta = 0.5$ and max-over-organs aggregation a priori. They are defensible but arbitrary, and no
sensitivity analysis over them is reported.

\textbf{Cohort.} We evaluate our index on 46 patients from a single collection,
with no external validation and no comparison against a clinical reference
standard for organ-involvement risk. Our per-level analysis rests on 7--17
lesions per level.

\textbf{Calibration provenance.} We originally fitted our per-organ constants on
the evaluation set, which made earlier versions of these numbers circular. We
now fit them on 35 held-out patients and every number here uses the corrected
constants, but the history is part of the record.

\textbf{Our expected-organ table may still be imperfect.} We found one error in
it and corrected it from anatomical literature, and that correction moved our
headline mean rank by 0.16 and shifted an entire level from below chance to
above it. The existence of one such error is evidence that others may remain, particularly in the same thoracolumbar region where the anatomy is
genuinely ambiguous. This is a real limit on how finely the per-level results
should be read.

\subsection{Future work}

The per-organ-scoped negative variant described above is the most specific next
step, because it addresses a mechanism we diagnosed rather than a
guess. Beyond that, we would harden our bilateral prompts with
position-invariant morphological cues, as our finding 3a implies, and adopt the
adversarial swap as a routine pre-deployment check. We would retrain with
minimal captions to resolve the prompt confound above. We would extend our index
to bone tumors as a distinct tissue target. And we would seek external
validation on an independent cohort, which the index most needs before any
clinical claim could be entertained.
